% Generated draft: requires mathematical and source review.
\documentclass[11pt]{article}
\usepackage[margin=25mm]{geometry}
\usepackage{fontspec}
\setmainfont{Latin Modern Roman}
\usepackage{amsmath,amssymb,mathtools,unicode-math}
\setmathfont{Latin Modern Math}
\usepackage{enumitem,needspace,xurl,hyperref}
\setlength{\parindent}{0pt}
\setlength{\parskip}{5pt}
\setlength{\emergencystretch}{3em}
\newcommand{\N}{\mathbb{N}}
\newcommand{\Z}{\mathbb{Z}}
\newcommand{\Q}{\mathbb{Q}}
\newcommand{\R}{\mathbb{R}}
\newcommand{\C}{\mathbb{C}}
\newcommand{\F}{\mathbb{F}}
\newcommand{\A}{\mathbb{A}}
\newcommand{\PP}{\mathbb{P}}
\newcommand{\E}{\mathbb{E}}
\newcommand{\eps}{\varepsilon}
\newcommand{\dd}{\,\mathrm{d}}
\begin{document}
\section*{UnsolvedMath ID: 1200024}
\textbf{Source problem code:} catalogue entry
\subsection{Definitions and mathematical statement}
For a connected locally finite simple graph $G$, the free uniform spanning forest $\mathfrak F_G$ is the weak limit of uniform spanning trees on finite connected subgraphs exhausting $G$, using free boundary conditions. Vertex transitivity means that for any two vertices an automorphism takes one to the other. Consequently
\[
\rho(G)=\mathbb E\big[\deg_{\mathfrak F_G}(v)\big]
\]
is independent of the choice of vertex $v$. The square $G^2$ has the same vertex set as $G$, with distinct vertices adjacent exactly when their $G$-distance is $1$ or $2$; it is not a Cartesian or direct graph product.

In the conventional infinite-graph forest normalization one writes
\[
\beta_1^{(2)}(G)=\tfrac12\rho(G)-1.
\]
The source asks for a definition from expected forest degree but does not supply this normalization, finite-graph conventions, or any trace/unimodularity hypotheses needed to identify the expression with a particular $L^2$-homology definition. Those choices must be fixed before claiming that identification.

\textbf{Question.} For every vertex-transitive graph in the intended class, is
\[
\beta_1^{(2)}(G)=\beta_1^{(2)}(G^2)?
\]
Under the displayed common forest normalization, the precise probabilistic target is
\[
\mathbb E\big[\deg_{\mathfrak F_G}(v)\big]
=\mathbb E\big[\deg_{\mathfrak F_{G^2}}(v)\big].
\]
The supplied summary reports no published resolution found within its search limit and describes the equality as apparently open. No conclusion is inferred merely from a comparison of the two graph Laplacians.
\subsection{Short English statement}
Does adding edges between vertices at distance two preserve the expected free-uniform-spanning-forest degree, and hence the first $L^2$ Betti number under the intended common normalization?
\subsection{Sources}
\begin{itemize}
\item M. Ab\'ert, \emph{Some questions} (2010), Question 24, \url{https://www.renyi.hu/~abert/questions.pdf}; manual transcription, accessed 2026-07-29. Current status was not independently verified in the source background.
\item The supplied research summary attributes the forest interpretation to Lyons but supplies no publication title, year, or theorem reference for its definition or scope.
\item Local registry \texttt{data/unsolvedmath/problems.json}, permanent ID $1200024$, source code catalogue entry, \emph{Some Questions --- Question 24}. The record labels this open and OPEN-TRIAGE, retaining source-status and rights qualifications marked NEEDS\_REVIEW.
\end{itemize}
\end{document}
